\section{Wnioski}\label{Wnioski}

Przedstawione poniżej wnioski końcowe stanowią bezpośrednie podsumowanie realizacji celu głównego oraz celów szczegółowych pracy, a także zawierają odpowiedź na pytanie badawcze i weryfikację postawionej hipotezy.

\subsection{Odpowiedź na pytanie badawcze}

Studenci przyjmują suplementy diety z bardzo wysoką częstotliwością (90,7\% próby, w tym blisko 65\% codziennie), traktując je jako stały element dbania o zdrowie. Choć większość z nich deklaruje wybrane zachowania bezpieczne (głównie analizę składu), to ogólny poziom stosowania się do procedur bezpieczeństwa jest niewystarczający – aż 41,24\% respondentów charakteryzuje się niską prawidłowością praktyk suplementacyjnych.

\subsection{Weryfikacja hipotezy}

Postawiona hipoteza badawcza, zakładająca, że wyższy poziom wiedzy studentów na temat suplementów diety wiąże się z bardziej racjonalnymi praktykami suplementacyjnymi, nie została potwierdzona. Analiza statystyczna wykazała brak istotnej zależności pomiędzy indeksem wiedzy a ogólnym indeksem prawidłowych praktyk ($\rho = 0,15$; $p = 0,11$). Wyższa wiedza sprzyja wprawdzie pojedynczym zachowaniom (sprawdzanie powielania składników), lecz nie przekłada się na całościowe, bezpieczne nawyki, takie jak diagnostyka laboratoryjna czy konsultacje lekarskie.

\subsection{Wnioski szczegółowe}

\begin{enumerate}
    \item Częstość stosowania (Cel 1): Suplementacja diety jest zjawiskiem powszechnym i utrwalonym w badanej grupie studenckiej. Ponad 90\% respondentów deklaruje aktualne przyjmowanie preparatów, z czego większość stosuje je codziennie oraz nieprzerwanie od ponad 3 lat. Zjawisko to ma charakter powszechny i nie zależy w sposób istotny statystycznie od płci, wieku czy poziomu aktywności fizycznej badanych.
    \item Struktura preferencji (Cel 2): Studenci najchętniej wybierają suplementy jednoskładnikowe i celowane o udowodnionym znaczeniu populacyjnym. Najpopularniejszymi preparatami są: witamina D (81,4\%), kwasy omega-3 (60,8\%) oraz magnez (49,4\%), co świadczy o profilu konsumenckim odchodzącym od tradycyjnych, wieloskładnikowych preparatów multiwitaminowych.
    \item Motywacje do suplementacji (Cel 3): Główną siłą napędową suplementacji są motywy prozdrowotne i profilaktyczne – poprawa zdrowia (70,1\%), wzmocnienie odporności (68,04\%) oraz zapobieganie niedoborom (67,012\%). Fakt, że wysoka suplementacja współwystępuje u części osób z niską aktywnością fizyczną, pozwala wnioskować, iż preparaty te bywają podświadomie traktowane jako łatwy substytut całościowego stylu życia.
    \item Źródła informacji (Cel 4): Internet i media społecznościowe stanowią dominujące źródło wiedzy o suplementach (łącznie blisko 70\% wskazań). Co istotne, deklarowane korzystanie z profesjonalnych źródeł (lekarz, dietetyk, farmaceuta) nie wiązało się w sposób istotny statystycznie z wyższym poziomem obiektywnej wiedzy w badanej próbie ($p = 0,29$).
    \item Poziom wiedzy studentów (Cel 5): Rzeczywisty poziom wiedzy studentów jest umiarkowany (średnia 3,76 na 6 punktów). Największą trudność sprawia badanym kwestia regulacji prawnych i procedur wprowadzania suplementów na rynek (tylko 50,12\% poprawnych odpowiedzi dotyczących braku wymogu badań klinicznych), co wskazuje na istnienie systemowej luki edukacyjnej.
    \item Wiedza a praktyka (Cel 6): Wyższy poziom wiedzy obiektywnej różnicuje badanych wyłącznie w zakresie sprawdzania powielania składników ($p = 0,0001$), nie wpływając istotnie na nawyki takie jak badania krwi ($p = 0,19$) czy konsultacje ze specjalistą ($p = 0,57$). Odnotowano natomiast silną, istotną statystycznie spójność między wiedzą rzeczywistą a subiektywną samooceną ($p < 0,001$) we wszystkich obszarach.
\end{enumerate}

\subsection{Rekomendacje aplikacyjne}

Wykryty w pracy dysonans (gdzie zadowalający poziom wiedzy teoretycznej i trafna samoocena nie gwarantują wdrażania bezpiecznych praktyk w codziennym życiu) implikuje wniosek, że przyszłe działania edukacyjne na uczelniach wyższych nie powinny ograniczać się do przekazywania suchych faktów o witaminach. Programy profilaktyczne muszą kłaść nacisk na behawioralne aspekty konsumpcji – zwalczanie poczucia złudnego bezpieczeństwa wobec produktów bezreceptowych oraz kształtowanie krytycznych postaw wobec marketingu. Jako preferowaną formę przekazu studenci wskazują krótkie formy wideo (54,64\%) oraz artykuły popularnonaukowe (46,39\%).

